\section{Zadanie 4}
Zmienna losowa \(X\) jest skoncentrowana na zbiorze \(D\), jeśli \(P(X \in D) = 1\). Mówimy, że zmienna losowa \(X\) skoncentrowana na \(\mathbb{N}\) jest bez pamięci jeśli dla każdego \(m, n \in N\) mamy

\[
	P(X > n) > 0 \text{ oraz } P(X = m + n |  X > n).
\]
Wykaż, że jeśli X jest bez pamięci to \(X\) ma rozkład geometycznych.

\subsection*{Rozwiązanie}
Naszym celem, jest pokazanie, że \(X\) ma rozkład geomentryczny. Dowiedziemy to, przez sprawdzenie ilorazu \(\frac{P(N+1)}{P(N)}\).

\[
	\frac{P(x=N+1)}{x=P(N)} = \frac{P(x=N+1)}{P(x=N+1 | x > 1)} =  \frac{P(x=N+1) \cdot P(x > 1)}{P(x = N+1 \wedge x > 1)} =
\]
\[
	\frac{P(x=N+1) \cdot P(x > 1)}{P(x=N+1)} = P(x > 1) = 1 - P(x = 1)
\]
Z tego wyszło, że kolejny element jest poprzednim przemonożonym przez \(1 - P(x = 1)\). Także z tego można wywniskować, że \(P(X = n) = (1-P(X=1))^{n-1} \cdot P(X = 1)\). Co jest definicją rozkładu geometrycznego.
